% ================================================================
%  БҮЛЭГ 1: СУДАЛГААНЫ ХЭСЭГ
% ================================================================
\chapter{СУДАЛГААНЫ ХЭСЭГ}
\label{ch1}

\section{Үндэслэл}

Хүнсний үйлчилгээний салбар нь дэлхийн нийт ДНБ-ний 3.4\%-ийг бүрдүүлж, 2023 оны
байдлаар 3.5 их наяд ам.долларын зах зээлийн хэмжээтэй байна~\cite{statista2023}.
Монгол улсад тухайн салбар ДНБ-ний ойролцоогоор 2.1\%-ийг бүрдүүлж байгаа бөгөөд
Улаанбаатар хотын 5,200 гаруй хүнсний үйлчилгээний газар 45,000 гаруй ажилтанг
хамарч байна~\cite{mrhba2023}.

COVID-19 цар тахлын үе нь дэлхийн рестораны салбарт гар хүрэлтгүй үйлчилгээний
эрэлтийг огцом нэмэгдүүлсэн бөгөөд QR кодоор тоног төхөөрөмж хоорондын харилцааг
хялбаршуулах боломж тодорхой болсон~\cite{gossling2022}. Азийн орнуудад QR кодын
ашиглалт 2020--2023 оны хооронд 185\%-иар өссөн бөгөөд ресторан, худалдаа, тээврийн
салбарт өргөн нэвтэрч байна~\cite{juniper2023}.

Монгол улсын рестораны 92\%-д шинэ технологийг нэвтрүүлэхэд санхүүгийн дарамт,
техникийн мэдлэгийн дутагдал, Монгол хэлийг дэмждэг бэлэн шийдлийн байхгүй байдал
зэрэг саад тотгор хэвээр байна. Иймд Монголын онцлогт тохирсон, хямд өртөгтэй,
нэвтрүүлэхэд хялбар шийдэл боловсруулах нь тулгамдсан асуудал болж байна.

Уламжлалт захиалгын арга нь дараах дөрвөн үндсэн асуудлыг бий болгодог:
\begin{itemize}
  \item \textbf{Захиалга алдагдах эрсдэл}: Зөөгч цаасан дэвтэрт бичиж авсан захиалга
    дутуу, эсвэл буруу дамжих нь ойролцоогоор 8\% байна~\cite{pos2022};
  \item \textbf{Хугацааны их зарцуулалт}: Нэг захиалга авахаас гал тогоонд очих хүртэлх
    хугацаа дунджаар 4.2 минут болдог~\cite{pos2022};
  \item \textbf{Хүний нөөцийн дарамт}: Нэг зөөгч дунджаар 8--10 ширээ хариуцдаг тул
    оргил цагт үйлчилгээний чанар буурдаг;
  \item \textbf{Олон хэлний хязгаарлалт}: Гадаадын зочдод цэсийг тайлбарлах хүнсний
    нэршлийн хэлний бэрхшээл байнга тулгардаг.
\end{itemize}

\section{Системийн зорилго}

Энэхүү дипломын ажлын гол зорилго нь рестораны үйлчлүүлэгч ширээн дэлгэц дээрх QR
кодыг смартфоноор уншуулахад тухайн рестораны дижитал цэс нэн даруй нээгдэж, хоол
сонгон захиалга өгөхөд тухайн захиалга нь бодит цаг хугацаанд ажилтны дэлгэцэнд
дамждаг бүрэн автоматжсан системийг загварчилж хэрэгжүүлэх явдал юм.

\section{Системийн зорилт}

\begin{enumerate}[leftmargin=*]
  \item Холбогдох шинжлэх ухааны бүтээлд дүн шинжилгээ хийж, санал болгах системийн
        онолын үндэслэлийг тогтоох;
  \item Хэрэглэгчийн болон техникийн шаардлагыг тодорхойлох;
  \item QR кодоос захиалга ажилтанд хүртэлх бүрэн урсгалын системийн загварыг
        боловсруулах;
  \item 4 давхаргат системийн архитектурыг хэрэгжүүлэх;
  \item Уламжлалт системтэй харьцуулан туршилтын үр дүнд дүн шинжилгээ хийх.
\end{enumerate}

\section{Ижил төстэй системийн судалгаа}

\subsection{Монголын системүүдтэй харьцуулсан судалгаа}

Монгол улсад хоол захиалгын системийн хувьд голлон уламжлалт арга буюу цайчинд
суурилсан захиалга зонхилж байна. Хэдийгээр Facebook, Instagram зэрэг сошиал медиагаар
онлайн хүргэлтийн захиалга авдаг рестораны тоо нэмэгдсэн ч ресторан доторх ширээнээс
шууд захиалга авдаг системийн хэрэгжилт хаана ч байдаггүй.

Нямбаяр нар~\cite{nyambayar2022} нарын 2022 оны судалгаагаар Улаанбаатар хотын
рестораны дижитал шилжилтийн бэлэн байдалд дүн шинжилгээ хийж, QR нэвтрүүлэх
суурь нөхцөл (смартфоны өндөр хүртээмж -- 78.4\%, 4G сүлжээний тогтвортой байдал)
бүрдсэн боловч тохиромжтой нутагшуулсан шийдэл байхгүй байгааг онцолжээ.

\begin{table}[htbp]
\caption{Монголын болон гадаад системүүдийн харьцуулалт}
\label{tab:compare}
\centering
\small
\begin{tabularx}{\textwidth}{|L|Y|Y|Y|Y|}
\hline
\rowcolor{TableHdr}
\textcolor{white}{\textbf{Шалгуур}} &
\textcolor{white}{\textbf{Toast POS}} &
\textcolor{white}{\textbf{Lightspeed}} &
\textcolor{white}{\textbf{Square}} &
\textcolor{white}{\textbf{Энэхүү систем}} \\ \hline
QR захиалга & Нэмэлт & Нэмэлт & Тийм & Гол функц \\ \hline
\rowcolor{LightBlue}
Бодит цаг WebSocket & Хэсэгчлэн & Тийм & Хэсэгчлэн & Тийм \\ \hline
Монгол хэл & Үгүй & Үгүй & Үгүй & Тийм \\ \hline
\rowcolor{LightBlue}
QPay / SocialPay & Үгүй & Үгүй & Үгүй & Тийм \\ \hline
Сарын зардал (USD) & 69--165 & 69--135 & 60+ & $\sim$10--25 \\ \hline
\rowcolor{LightBlue}
Нээлттэй эх код & Үгүй & Үгүй & Үгүй & Тийм \\ \hline
\end{tabularx}
\end{table}

\subsection{Гадны системүүдтэй харьцуулсан судалгаа}

\textbf{Toast POS (АНУ, 2012)} нь АНУ-ын өргөн хэрэглэгддэг POS системүүдийн нэг.
QR кодоор захиалга авах боломжтой боловч нэмэлт модул шаарддаг. Монгол хэл болон
QPay, SocialPay зэрэг локал төлбөрийн хэрэгсэл байхгүй~\cite{toast2023}.

\textbf{Lightspeed Restaurant (Канад, 2005)} нь 115,000 гаруй рестораны газарт
хэрэглэгддэг. Сарын суурь хөлс 69 ам.доллар бөгөөд QR захиалгын модул нэмэлтийн
39 ам.долларын хөлстэй. Гадаад валютаар тооцогдох зардлын дарамт практикт
саад учруулна~\cite{lightspeed2023}.

\textbf{Square for Restaurants (АНУ, 2018)} нь QR захиалгыг дэмждэг боловч гадаад
картын системд суурилсан тул Монгол улсын картын дэд бүтцэд зориулагдаагүй~\cite{square2023}.

\textbf{Ajisen Ramen (Хятад, 2019)} нь Хятадад амжилттай хэрэгжсэн QR захиалгын
тогтолцоотой. Хэдий тийм ч систем нь WeChat болон Alipay-тэй нягт холбоотой тул
Монголд шилжүүлэх боломжгүй~\cite{ajisen2020}.

\section{Хууль дүрэм журам}

\subsection{Хувийн мэдээлэл хамгаалах тухай хууль}

Монгол Улсын ``Хувийн мэдээлэл хамгаалах тухай'' хуулийн (2021) 14.1 дүгээр зүйлд
иргэний хувийн мэдээллийг эзэмшигчийн зөвшөөрөлгүй хадгалах, дамжуулахыг
хориглосон байна~\cite{pil2021}. Тус систем хэрэглэгчийн нэр, утасны дугаар, байршлын
мэдээллийг хадгалахгүй бөгөөд зөвхөн тухайн ширээний QR token-оор захиалгыг
таних тул хуулийн шаардлагыг бүрэн хангана.

\subsection{Кибер аюулгүй байдлын тухай хууль}

``Кибер аюулгүй байдлын тухай'' хуулийн (2021) 7.1 дүгээр зүйлд мэдээллийн систем
нь нэвтрэх нотлох, шифрлэлт, бүртгэлийн механизмтай байхыг шаарддаг~\cite{cslaw2021}.
Тус систем JWT HS256 шифрлэлт, HTTPS/TLS, хандалтын бүртгэл (audit log) зэрэг
хамгаалалтыг нэвтрүүлсэн.

\subsection{Цахим төлбөр тооцооны тухай хууль}

``Цахим төлбөр тооцооны тухай'' хуулийн (2017) 8.2 дугаар зүйлд цахим гүйлгээний
бүртгэлийг таван жилийн хугацаанд хадгалахыг шаарддаг~\cite{eplaw2017}. Системийн
бүх захиалга, төлбөрийн мэдээлэл өгөгдлийн санд архивлагдан хадгалагдана.

\subsection{Хүнсний тухай хууль}

``Хүнсний тухай'' хуульд (2012) хүнсний бүтээгдэхүүний найрлага, хадгалах нөхцлийг
дижитал цэсэд тусгахыг зөвлөдөг~\cite{foodlaw2012}. Тус систем хоол бүрийн дэлгэрэнгүй
тайлбар, аллерген мэдээлэл оруулах боломжийг хэрэгжүүлсэн.

\subsection{Хэрэглэгчийн эрхийг хамгаалах тухай хууль}

``Хэрэглэгчийн эрхийг хамгаалах тухай'' хуулийн (2003) 8.1 дүгээр зүйлд үнийн
ил тод байдлыг шаарддаг~\cite{cpal2003}. Тус систем цэс дэх бүх бүтээгдэхүүний
үнийг нэн даруй харуулж, нийт дүнг захиалахаас өмнө баталгаажуулдаг.

\section{Онолын хэсэг}

\subsection{QR кодын онол ба технологи}

QR (Quick Response) код нь 1994 онд Японы Denso Wave компаний Масахиро Хара боловсруулсан
2 хэмжээст матрицан штрих код юм~\cite{iso18004}. QR кодын үндсэн онцлог:

\begin{itemize}
  \item \textbf{Мэдээллийн нягтрал}: Нэг QR кодод хамгийн ихдээ 4,296 тэмдэгт
    (alphanumeric), 7,089 цифр агуулах боломжтой;
  \item \textbf{Алдааны засалт}: Reed-Solomon алгоритм дээр суурилан 7\%--30\% алдааны
    засалтын түвшинтэй (L, M, Q, H гэсэн 4 түвшин);
  \item \textbf{Чиглэлийн мэдрэлт}: Гурван булан дах тогтсон загварт дурын чиглэлийн
    мэдрэлт хийгдэнэ;
  \item \textbf{Нэвтрэх хурд}: ISO 18004 стандартад тодорхойлсноор 0.5 секундын дотор
    задлагдана~\cite{iso18004}.
\end{itemize}

Рестораны систем дэх QR кодын URL бүтэц дараах байдлаар зохиомжлогдсон:

\begin{center}
\texttt{https://app.restaurant.mn/menu?t=\{qr\_token\}}
\end{center}

Энд \texttt{qr\_token} нь UUID форматтай, ширээ бүрт өвөрмөц утгатай, серверийн
өгөгдлийн санд хадгалагдах тоон утга юм.

\subsection{WebSocket ба бодит цагийн технологи}

WebSocket нь 2011 онд IETF-ийн RFC 6455 стандартаар батлагдсан, HTTP-ийн дээр
суурилсан 2 талын байнгын холболтын протокол юм~\cite{rfc6455}. Уламжлалт HTTP
``хүсэлт-хариу'' загвараас ялгаатай нь WebSocket холболт нэгэнт тогтсоны дараа
сервер өөрөө клиентэд мессеж илгээх боломжтой.

\begin{figure}[htbp]
\centering
\begin{tikzpicture}[
  box/.style={rectangle, draw, rounded corners=4pt, minimum width=3cm,
    minimum height=0.8cm, align=center, font=\small},
  arr/.style={-{Stealth}, thick}
]
\node[box, fill=blue!15] (client) at (0, 0) {Клиент (Browser)};
\node[box, fill=green!15] (server) at (8, 0) {Сервер (Socket.IO)};

\draw[arr] (client) to[bend left=20] node[above,font=\scriptsize]
  {HTTP Upgrade Request} (server);
\draw[arr] (server) to[bend left=20] node[below,font=\scriptsize]
  {101 Switching Protocols} (client);

\draw[arr] (client) -- node[above,font=\scriptsize,pos=0.4]
  {emit('place\_order')} ($(client)!0.5!(server)+(0,-1)$);
\draw[arr] ($(client)!0.5!(server)+(0,-1)$) -- (server);

\draw[arr] (server) to[bend left=15] node[above,font=\scriptsize]
  {emit('order\_update')} (client);

\node[font=\tiny,text=gray] at (4,-2) {TCP байнгын холболт -- header overhead байхгүй};
\end{tikzpicture}
\caption{WebSocket протоколын харилцааны загвар}
\label{fig:websocket}
\end{figure}

Socket.IO нь WebSocket-ийн дээр суурилан дараах давхар боломжуудыг нэмдэг:
автомат дахин холболт, өрөө (room)-д суурилсан дохио, HTTP long-polling буцаах
(fallback), namespace дэмжлэг~\cite{socketio2023}.

\subsection{REST API архитектурын онол}

REST (Representational State Transfer) нь Рой Филдингийн 2000 оны докторын
ажилд тодорхойлсон вэб үйлчилгээний архитектурын загвар~\cite{fielding2000}.
REST-ийн 6 үндсэн зарчим:

\begin{enumerate}[leftmargin=*]
  \item \textbf{Client-Server}: Клиент болон серверийг тусгаарлах нь бие даасан
    хөгжлийг боломжтой болгодог;
  \item \textbf{Stateless}: Хүсэлт бүр бие даасан, серверт session хадгалагддаггүй;
  \item \textbf{Cacheable}: Хариу мэдэгдлийг кэш хийх боломжтой;
  \item \textbf{Uniform Interface}: HTTP методоор (GET, POST, PUT, PATCH, DELETE)
    нэгдмэл интерфейс;
  \item \textbf{Layered System}: Дэд бүтцийн давхаргууд клиентэд харагдахгүй;
  \item \textbf{Code on Demand}: Сервер скрипт илгээх боломжтой (заавал биш).
\end{enumerate}

\section{Технологийн судалгаа}

Системийг хөгжүүлэх технологийг сонгохдоо Монголын хөгжүүлэгчдийн экосистем,
нийгэмлэгийн дэмжлэг, чөлөөт тусламжийн боломж, хямд өртөг зэргийг харгалзсан.

\begin{table}[htbp]
\caption{Ашигласан технологиудын харьцуулалт ба сонголтын үндэслэл}
\label{tab:tech}
\centering
\small
\begin{tabularx}{\textwidth}{|L|L|Y|}
\hline
\rowcolor{TableHdr}
\textcolor{white}{\textbf{Давхарга}} &
\textcolor{white}{\textbf{Сонгосон технологи}} &
\textcolor{white}{\textbf{Үндэслэл}} \\ \hline
Хэрэглэгчийн интерфейс & React 18 + TypeScript &
  Component-based, virtual DOM, өргөн экосистем \\ \hline
\rowcolor{LightBlue}
Хэрэглэгчийн интерфейс & Tailwind CSS v4 &
  Utility-first, хурдан загвар боловсруулалт \\ \hline
API сервер & Node.js + Express.js &
  JavaScript нэгдсэн стек, event-driven I/O \\ \hline
\rowcolor{LightBlue}
Бодит цаг & Socket.IO v4 &
  WebSocket + fallback, room дэмжлэг \\ \hline
Өгөгдлийн сан & PostgreSQL 16 &
  ACID, JSON дэмжлэг, хүчирхэг хувилбарын хяналт \\ \hline
\rowcolor{LightBlue}
ORM & Drizzle ORM &
  TypeScript native, migration автомат \\ \hline
Баталгаажуулалт & JWT HS256 &
  Stateless, хурдан шалгалт \\ \hline
\rowcolor{LightBlue}
Багц удирдлага & pnpm workspace &
  Монорепо, хуваалцах dependency \\ \hline
\end{tabularx}
\end{table}

\textbf{Frontend технологийн сонголт:} React 18-ийн Concurrent Mode болон Suspense
механизм нь динамик цэс ачааллах хугацааг 40\%-иар бууруулдаг. TypeScript-ийн статик
шалгалт нь хөгжүүлэлтийн алдааны тоог 23\%-иар бууруулдаг талаар Stack Overflow-ийн
2023 оны судалгаанд дурдсан байдаг~\cite{stackoverflow2023}.

\textbf{Backend технологийн сонголт:} Node.js-ийн libuv event loop нь хавтгай I/O
дарааллаар олон нэгэн цагт хүсэлтийг шийдвэрлэдэг. PostgreSQL 16-ийн Logical
Replication болон JSON дэмжлэг нь цаашдын scale out-ийн суурь болно.

\textbf{Socket.IO-ийн давуу тал:} Нэг Socket.IO үйлдлийн сервер 10,000 нэгэн цагт
холболтыг дэмждэг бөгөөд монгол рестораны хэмжээнд (дунджаар 30--100 ширээ)
нэмэлт дэд бүтцийн хөрөнгө оруулалтгүйгээр ажиллана~\cite{socketio2023}.

\section{Бүлгийн дүгнэлт}

Энэхүү бүлэгт судалгааны үндэслэл, зорилго, зорилтыг тодорхойлж, одоо байгаа
гадаад болон дотоодын ижил зориулалтын системүүдтэй харьцуулан дүгнэлт хийв.
Монгол Улсын холбогдох хууль тогтоомжийн шаардлагыг нарийвчлан судалж, системийн
хэрэгжилт тэдгээрийг бүрэн хангаж байгааг тогтоов. WebSocket болон REST API-ийн
онолын үндэслэлийг тайлбарлаж, ашигласан технологиудыг харьцуулан судалсны үндсэн
дээр сонголт хийсэн болно.
